
\documentclass{article}
\usepackage{colortbl}
\usepackage{makecell}
\usepackage{multirow}
\usepackage{supertabular}

\begin{document}

\newcounter{utterance}

\centering \large Interaction Transcript for game `dond', experiment `semi', episode 0 with Qwen3.6{-}35B{-}A3B{-}FP8{-}without{-}reasoning.
\vspace{24pt}

{ \footnotesize  \setcounter{utterance}{1}
\setlength{\tabcolsep}{0pt}
\begin{supertabular}{c@{$\;$}|p{.15\linewidth}@{}p{.15\linewidth}p{.15\linewidth}p{.15\linewidth}p{.15\linewidth}p{.15\linewidth}}
    \# & $\;$A & \multicolumn{4}{c}{Game Master} & $\;\:$B\\
    \hline

    \theutterance \stepcounter{utterance}  
    & & \multicolumn{4}{p{0.6\linewidth}}{
        \cellcolor[rgb]{0.9,0.9,0.9}{
            \makecell[{{p{\linewidth}}}]{
                \texttt{\tiny{[P1$\langle$GM]}}
                \texttt{Ви граєте в кооперативну переговорну гру, у якій вам потрібно домовитися, як поділити набір предметів між вами та іншим гравцем.} \\
\\ 
\texttt{Правила:} \\
\texttt{(a) Вам та іншому гравцеві надається набір предметів. Кожному з вас також надається секретна функція цінності, що відображає, наскільки ви цінуєте кожен тип предмета.} \\
\texttt{(b) Ви обмінюєтеся повідомленнями з іншим гравцем, щоб домовитися, кому які предмети дістануться. Кожен із вас може надіслати максимум 5 повідомлень або завершити раніше, зробивши секретну пропозицію в будь{-}який момент.} \\
\texttt{(c) Кожного з вас просять подати секретну пропозицію із зазначенням предметів, які ви хочете отримати, у форматі в квадратних дужках, як показано нижче: "{[}Пропозиція: <число> <назва об'єкта>, <число> <назва об'єкта>, <...>{]}"} \\
\texttt{(d) Якщо ваші пропозиції є взаємодоповнювальними, тобто предметів достатньо, щоб виконати обидві пропозиції, кожному гравцеві нараховуються бали на основі суми цінностей предметів, які він отримав. Інакше ви обоє отримуєте нуль балів.} \\
\texttt{(e) Ваша мета — максимізувати суму вашого рахунку та рахунку іншого гравця.} \\
\\ 
\texttt{Почнімо.} \\
\\ 
\texttt{Набір доступних предметів:} \\
\\ 
\texttt{3 кільця, 1 кубик, 1 сковорода} \\
\\ 
\texttt{Ваша секретна функція цінності:} \\
\\ 
\texttt{каблучка: 0, кубик: 4, сковорода: 6} \\
\\ 
\texttt{ВАЖЛИВО: Ваші повідомлення, якщо це не секретна пропозиція, безпосередньо передаються іншому гравцеві, тож не включайте жодної відповіді на правила або тексту, що оголошує ваше повідомлення. Щоб зробити секретну пропозицію, використовуйте вказаний формат. Не використовуйте квадратні дужки під час спілкування з іншим гравцем, інакше це буде інтерпретовано як ваша секретна пропозиція.} \\
            }
        }
    }
    & \\ \\

    \theutterance \stepcounter{utterance}  
    & \multicolumn{4}{p{0.6\linewidth}}{
        \cellcolor[rgb]{0.9,0.9,0.9}{
            \makecell[{{p{\linewidth}}}]{
                \texttt{\tiny{[P1$\rangle$GM]}}
                \texttt{Привіт! Я бачу, що у нас є 3 кільця, 1 кубик і 1 сковорода.} \\
\\ 
\texttt{Згідно з моїми оцінками, сковорода для мене найцінніша (6 балів), а кубик — 4 бали. Кільця я ціную дуже низько (0 балів).} \\
\\ 
\texttt{Щоб ми обидва отримали максимальну користь, я пропоную домовитися так:} \\
\texttt{1. Я беру сковороду.} \\
\texttt{2. Ти можеш взяти кубик.} \\
\texttt{3. Кільця можуть дістатися тобі, оскільки для мене вони не мають цінності, а для тебе, сподіваюся, вони цінніші.} \\
\\ 
\texttt{Якщо ти згоден з таким розподілом (я беру сковороду, ти береш кубик і кільця), будь ласка, підтвердь це або запропонуй свій варіант. Я готовий надіслати секретну пропозицію, коли ми домовимося.} \\
            }
        }
    }
    & & \\ \\

    \theutterance \stepcounter{utterance}  
    & & & \multicolumn{2}{p{0.3\linewidth}}{
        \cellcolor[rgb]{0.9,0.9,0.9}{
            \makecell[{{p{\linewidth}}}]{
                \texttt{\tiny{[GM$|$GM]}}
                \texttt{continue} \\
            }
        }
    }
    & & \\ \\

    \theutterance \stepcounter{utterance}  
    & & & \multicolumn{2}{p{0.3\linewidth}}{
        \cellcolor[rgb]{0.9,0.9,0.9}{
            \makecell[{{p{\linewidth}}}]{
                \texttt{\tiny{[GM$|$GM]}}
                \texttt{Привіт! Я бачу, що у нас є 3 кільця, 1 кубик і 1 сковорода.} \\
\\ 
\texttt{Згідно з моїми оцінками, сковорода для мене найцінніша (6 балів), а кубик — 4 бали. Кільця я ціную дуже низько (0 балів).} \\
\\ 
\texttt{Щоб ми обидва отримали максимальну користь, я пропоную домовитися так:} \\
\texttt{1. Я беру сковороду.} \\
\texttt{2. Ти можеш взяти кубик.} \\
\texttt{3. Кільця можуть дістатися тобі, оскільки для мене вони не мають цінності, а для тебе, сподіваюся, вони цінніші.} \\
\\ 
\texttt{Якщо ти згоден з таким розподілом (я беру сковороду, ти береш кубик і кільця), будь ласка, підтвердь це або запропонуй свій варіант. Я готовий надіслати секретну пропозицію, коли ми домовимося.} \\
            }
        }
    }
    & & \\ \\

    \theutterance \stepcounter{utterance}  
    & & \multicolumn{4}{p{0.6\linewidth}}{
        \cellcolor[rgb]{0.9,0.9,0.9}{
            \makecell[{{p{\linewidth}}}]{
                \texttt{\tiny{[GM$\rangle$P2]}}
                \texttt{Ви граєте в кооперативну переговорну гру, у якій вам потрібно домовитися, як поділити набір предметів між вами та іншим гравцем.} \\
\\ 
\texttt{Правила:} \\
\texttt{(a) Вам та іншому гравцеві надається набір предметів. Кожному з вас також надається секретна функція цінності, що відображає, наскільки ви цінуєте кожен тип предмета.} \\
\texttt{(b) Ви обмінюєтеся повідомленнями з іншим гравцем, щоб домовитися, кому які предмети дістануться. Кожен із вас може надіслати максимум 5 повідомлень або завершити раніше, зробивши секретну пропозицію в будь{-}який момент.} \\
\texttt{(c) Кожного з вас просять подати секретну пропозицію із зазначенням предметів, які ви хочете отримати, у форматі в квадратних дужках, як показано нижче: "{[}Пропозиція: <число> <назва об'єкта>, <число> <назва об'єкта>, <...>{]}"} \\
\texttt{(d) Якщо ваші пропозиції є взаємодоповнювальними, тобто предметів достатньо, щоб виконати обидві пропозиції, кожному гравцеві нараховуються бали на основі суми цінностей предметів, які він отримав. Інакше ви обоє отримуєте нуль балів.} \\
\texttt{(e) Ваша мета — максимізувати суму вашого рахунку та рахунку іншого гравця.} \\
\\ 
\texttt{Почнімо.} \\
\\ 
\texttt{Набір доступних предметів:} \\
\\ 
\texttt{3 кільця, 1 кубик, 1 сковорода} \\
\\ 
\texttt{Ваша секретна функція цінності:} \\
\\ 
\texttt{каблучка: 3, кубик: 1, сковорода: 0} \\
\\ 
\texttt{ВАЖЛИВО: Ваші повідомлення, якщо це не секретна пропозиція, безпосередньо передаються іншому гравцеві, тож не включайте жодної відповіді на правила або тексту, що оголошує ваше повідомлення. Щоб зробити секретну пропозицію, використовуйте вказаний формат. Не використовуйте квадратні дужки під час спілкування з іншим гравцем, інакше це буде інтерпретовано як ваша секретна пропозиція.} \\
\\ 
\texttt{Привіт! Я бачу, що у нас є 3 кільця, 1 кубик і 1 сковорода.} \\
\\ 
\texttt{Згідно з моїми оцінками, сковорода для мене найцінніша (6 балів), а кубик — 4 бали. Кільця я ціную дуже низько (0 балів).} \\
\\ 
\texttt{Щоб ми обидва отримали максимальну користь, я пропоную домовитися так:} \\
\texttt{1. Я беру сковороду.} \\
\texttt{2. Ти можеш взяти кубик.} \\
\texttt{3. Кільця можуть дістатися тобі, оскільки для мене вони не мають цінності, а для тебе, сподіваюся, вони цінніші.} \\
\\ 
\texttt{Якщо ти згоден з таким розподілом (я беру сковороду, ти береш кубик і кільця), будь ласка, підтвердь це або запропонуй свій варіант. Я готовий надіслати секретну пропозицію, коли ми домовимося.} \\
            }
        }
    }
    & \\ \\

    \theutterance \stepcounter{utterance}  
    & & & \multicolumn{4}{p{0.6\linewidth}}{
        \cellcolor[rgb]{0.9,0.9,0.9}{
            \makecell[{{p{\linewidth}}}]{
                \texttt{\tiny{[GM$\langle$P2]}}
                \texttt{Привіт! Дякую за пропозицію.} \\
\\ 
\texttt{Однак я хочу уточнити свої оцінки, щоб ми могли знайти найкращий спосіб розподілу. Для мене найціннішим є кубик (5 балів), кільця мають середню цінність (3 бали кожне), а сковорода для мене не має цінності (0 балів).} \\
\\ 
\texttt{Тобі, схоже, найбільше потрібна сковорода. Якщо ти готовий віддати мені кубик, я можу віддати тобі всі три кільця. Сковороду тоді залишаємо тобі.} \\
\\ 
\texttt{Якщо ти згоден на такий варіант (я беру кубик, ти береш сковороду та 3 кільця), будь ласка, підтвердь, і я надішлю секретну пропозицію.} \\
            }
        }
    }
     \\ \\

    \theutterance \stepcounter{utterance}  
    & & & \multicolumn{2}{p{0.3\linewidth}}{
        \cellcolor[rgb]{0.9,0.9,0.9}{
            \makecell[{{p{\linewidth}}}]{
                \texttt{\tiny{[GM$|$GM]}}
                \texttt{continue} \\
            }
        }
    }
    & & \\ \\

    \theutterance \stepcounter{utterance}  
    & & & \multicolumn{2}{p{0.3\linewidth}}{
        \cellcolor[rgb]{0.9,0.9,0.9}{
            \makecell[{{p{\linewidth}}}]{
                \texttt{\tiny{[GM$|$GM]}}
                \texttt{Привіт! Дякую за пропозицію.} \\
\\ 
\texttt{Однак я хочу уточнити свої оцінки, щоб ми могли знайти найкращий спосіб розподілу. Для мене найціннішим є кубик (5 балів), кільця мають середню цінність (3 бали кожне), а сковорода для мене не має цінності (0 балів).} \\
\\ 
\texttt{Тобі, схоже, найбільше потрібна сковорода. Якщо ти готовий віддати мені кубик, я можу віддати тобі всі три кільця. Сковороду тоді залишаємо тобі.} \\
\\ 
\texttt{Якщо ти згоден на такий варіант (я беру кубик, ти береш сковороду та 3 кільця), будь ласка, підтвердь, і я надішлю секретну пропозицію.} \\
            }
        }
    }
    & & \\ \\

    \theutterance \stepcounter{utterance}  
    & & \multicolumn{4}{p{0.6\linewidth}}{
        \cellcolor[rgb]{0.9,0.9,0.9}{
            \makecell[{{p{\linewidth}}}]{
                \texttt{\tiny{[P1$\langle$GM]}}
                \texttt{Привіт! Дякую за пропозицію.} \\
\\ 
\texttt{Однак я хочу уточнити свої оцінки, щоб ми могли знайти найкращий спосіб розподілу. Для мене найціннішим є кубик (5 балів), кільця мають середню цінність (3 бали кожне), а сковорода для мене не має цінності (0 балів).} \\
\\ 
\texttt{Тобі, схоже, найбільше потрібна сковорода. Якщо ти готовий віддати мені кубик, я можу віддати тобі всі три кільця. Сковороду тоді залишаємо тобі.} \\
\\ 
\texttt{Якщо ти згоден на такий варіант (я беру кубик, ти береш сковороду та 3 кільця), будь ласка, підтвердь, і я надішлю секретну пропозицію.} \\
            }
        }
    }
    & \\ \\

    \theutterance \stepcounter{utterance}  
    & \multicolumn{4}{p{0.6\linewidth}}{
        \cellcolor[rgb]{0.9,0.9,0.9}{
            \makecell[{{p{\linewidth}}}]{
                \texttt{\tiny{[P1$\rangle$GM]}}
                \texttt{Звучить чудово! Цей варіант ідеально підходить для обох нас.} \\
\\ 
\texttt{Я підтверджую: я беру сковороду, а ти береш кубик і всі три кільця.} \\
\\ 
\texttt{Ось моя секретна пропозиція:} \\
\\ 
\texttt{{[}Пропозиція: 1 сковорода{]}} \\
            }
        }
    }
    & & \\ \\

    \theutterance \stepcounter{utterance}  
    & & & \multicolumn{2}{p{0.3\linewidth}}{
        \cellcolor[rgb]{0.9,0.9,0.9}{
            \makecell[{{p{\linewidth}}}]{
                \texttt{\tiny{[GM$|$GM]}}
                \texttt{proposal} \\
            }
        }
    }
    & & \\ \\

    \theutterance \stepcounter{utterance}  
    & & & \multicolumn{2}{p{0.3\linewidth}}{
        \cellcolor[rgb]{0.9,0.9,0.9}{
            \makecell[{{p{\linewidth}}}]{
                \texttt{\tiny{[GM$|$GM]}}
                [0, 0, 1]
            }
        }
    }
    & & \\ \\

    \theutterance \stepcounter{utterance}  
    & & \multicolumn{4}{p{0.6\linewidth}}{
        \cellcolor[rgb]{0.9,0.9,0.9}{
            \makecell[{{p{\linewidth}}}]{
                \texttt{\tiny{[GM$\rangle$P2]}}
                \texttt{Інший гравець подав таємну пропозицію. Тепер ваша черга подати власну пропозицію.} \\
            }
        }
    }
    & \\ \\

    \theutterance \stepcounter{utterance}  
    & & & \multicolumn{4}{p{0.6\linewidth}}{
        \cellcolor[rgb]{0.9,0.9,0.9}{
            \makecell[{{p{\linewidth}}}]{
                \texttt{\tiny{[GM$\langle$P2]}}
                \texttt{{[}Пропозиція: 1 кубик{]}} \\
            }
        }
    }
     \\ \\

    \theutterance \stepcounter{utterance}  
    & & & \multicolumn{2}{p{0.3\linewidth}}{
        \cellcolor[rgb]{0.9,0.9,0.9}{
            \makecell[{{p{\linewidth}}}]{
                \texttt{\tiny{[GM$|$GM]}}
                \texttt{proposal} \\
            }
        }
    }
    & & \\ \\

    \theutterance \stepcounter{utterance}  
    & & & \multicolumn{2}{p{0.3\linewidth}}{
        \cellcolor[rgb]{0.9,0.9,0.9}{
            \makecell[{{p{\linewidth}}}]{
                \texttt{\tiny{[GM$|$GM]}}
                [0, 1, 0]
            }
        }
    }
    & & \\ \\

    \theutterance \stepcounter{utterance}  
    & & & \multicolumn{2}{p{0.3\linewidth}}{
        \cellcolor[rgb]{0.9,0.9,0.9}{
            \makecell[{{p{\linewidth}}}]{
                \texttt{\tiny{[GM$|$GM]}}
                [[0, 0, 1], [0, 1, 0]]
            }
        }
    }
    & & \\ \\

\end{supertabular}
}

\end{document}
